\section{Environment and Orchestration}
\label{sec:environment}
The experimental environment is strictly local and fully containerized, orchestrated via \texttt{docker compose} to ensure consistency and reproducibility across different host machines. The stack is segmented into profiles:
\begin{itemize}
    \item \texttt{--profile core}: Starts the foundational components, including the Hardhat L1 node, the deployer script, Executor, Submitter, and Sequencer.
    \item \texttt{--profile bench}: Triggers the Workload Generator and benchmark harness (\texttt{run\_experiment.sh}).
    \item \texttt{--profile report}: Triggers the Data Tools analytics pipeline to process raw telemetry into aggregated reports.
\end{itemize}

\subsection{Local L1 Configuration}
The local L1 network is provided by a custom Docker image (\texttt{contracts/Dockerfile.node}) running \texttt{npx hardhat node -{}-hostname 0.0.0.0 -{}-port 8545}. This setup is crucial because it provides deterministic, ultra-fast block mining (12-second intervals by default) and pre-funded accounts, allowing the benchmark to saturate the sequencer without being constrained by public testnet block times or faucet limits. Contract deployment is automated via a deployment script (\texttt{npx hardhat run scripts/deploy-local.ts}) executed immediately after the node initializes.

When running the local sequencer against a fresh local Hardhat node, the \texttt{start\_block} in \texttt{sequencer.yaml} (or \texttt{sequencer.docker.yaml}) must be strictly set to \texttt{0} to prevent the sequencer from attempting to sync historical blocks that do not exist, which would otherwise cause initialization failures.